(1) 将状态方程展开至 \(V^{-2}\)，有 \(P=RT/V+(RTb-a)/V^2+O(V^{-3})\)。倒转为低压形式 \(V=RT/P+b-a/(RT)+O(P)\)，故第二维里体积系数 \(B(T)=b-a/(RT)\)；它的单位必须是体积每摩尔。
(2) 对一摩尔体系，热力学恒等式给 \(dH=C_PdT+[V-T(\partial V/\partial T)_P]dP\)。令节流近似等焓 \(dH=0\)，便得 \(\mu_{\rm JT}=(\partial T/\partial P)_H=[T(\partial V/\partial T)_P-V]/C_P\)。正的 \(\mu_{\rm JT}\) 表示减压时 \(dT<0\)，即气体冷却。
(3) 由 \(V=RT/P+B(T)\)，理想气体部分在 \(T\partial_TV-V\) 中抵消，只剩 \(T B'(T)-B(T)=2a/(RT)-b\)。于是 \(\mu_{\rm JT}\simeq[2a/(RT)-b]/C_P\)；令它为零得低压反转温度 \(T_{\rm inv}=2a/(Rb)\)。在 \(T<T_{\rm inv}\) 时降压冷却，在较高温度时降压升温。
(4) 仅有排斥体积 \(a=0\) 时 \(\mu_{\rm JT}\simeq-b/C_P<0\)；仅有吸引 \(b=0\) 时 \(\mu_{\rm JT}\simeq2a/(RTC_P)>0\)，两者符号合理。推导假定 \(P\) 足够小、\(C_P\) 可用该态值且远离相变；高压非理想项、液化与临界涨落会改变反转曲线。
